\chapter{Quantum Observable Universe}
\label{chap:qou}

% ── Introduction ──────────────────────────────────────────────────────────────
\section{Introduction}
\label{sec:qou-intro}

The \emph{Quantum Observable Universe} (QOU) is the physical universe as
accessible through quantum-mechanical observation.  In this chapter we lay the
mathematical foundations for studying the QOU, including the relevant Hilbert
space structures, operator algebras, and the role of measurement.

% ── Foundations ───────────────────────────────────────────────────────────────
\section{Mathematical Foundations}
\label{sec:qou-foundations}

\begin{definition}[Observable Universe Hilbert Space]
  Let $\mathcal{H}_{\mathrm{QOU}}$ denote the separable Hilbert space of all
  square-integrable state vectors describing the quantum state of the
  observable universe.
\end{definition}

\begin{definition}[Quantum Observable]
  A \emph{quantum observable} is a self-adjoint operator
  $\hat{O}: \mathcal{H}_{\mathrm{QOU}} \to \mathcal{H}_{\mathrm{QOU}}$
  whose spectral decomposition encodes the possible measurement outcomes.
\end{definition}

% ── Structure ─────────────────────────────────────────────────────────────────
\section{Structure of the Observable Universe}
\label{sec:qou-structure}

% TODO: develop the algebraic structure of observables, including the
%       C*-algebra of observables and its representations.

\begin{remark}
  The algebra of observables $\mathcal{A}$ forms a $C^*$-algebra under the
  operator norm.  States on $\mathcal{A}$ are positive linear functionals of
  norm one; pure states correspond to rays in $\mathcal{H}_{\mathrm{QOU}}$.
\end{remark}

% ── Dynamics ──────────────────────────────────────────────────────────────────
\section{Dynamics}
\label{sec:qou-dynamics}

The time evolution of a state $\ket{\psi(t)} \in \mathcal{H}_{\mathrm{QOU}}$
is governed by the Schr\"{o}dinger equation:
\begin{equation}
  i\hbar \frac{\partial}{\partial t}\ket{\psi(t)}
  = \hat{H}\ket{\psi(t)},
  \label{eq:schrodinger}
\end{equation}
where $\hat{H}$ is the Hamiltonian of the universe.

% ── Summary ───────────────────────────────────────────────────────────────────
\section{Summary}
\label{sec:qou-summary}

% TODO: summarize key results and set up the connection to quantum torsion
%       discussed in Chapter~\ref{chap:torsion}.
